% ============================================================
%  Preambulo — Informe IEEE (dos columnas) · Proyecto final EL-5409 (Ball on Plate)
% ============================================================

\usepackage[T1]{fontenc}
\usepackage[utf8]{inputenc}
\usepackage[spanish,es-tabla,es-noquoting]{babel}
\usepackage{amsmath,amssymb}
\usepackage{graphicx}
\usepackage{booktabs}
\usepackage{array}
\usepackage{enumitem}
\usepackage{listings}
\usepackage{xcolor}
\usepackage{balance}
\usepackage{xurl}          % permite cortar URLs largas en la bibliografia
\usepackage[hidelinks,unicode]{hyperref}

% IEEEtran define \IEEEproof etc.; evitamos choques de babel
\addto\captionsspanish{\renewcommand{\tablename}{TABLA}}

% ---------- Unidades ----------
\usepackage{siunitx}
\sisetup{
  output-decimal-marker = {,},
  group-separator       = {\,},
  per-mode              = symbol,
  detect-all            = true
}

% ---------- Graficas ----------
\usepackage{pgfplots}
\usepackage{pgfplotstable}
\pgfplotsset{compat=1.18}
\usepgfplotslibrary{groupplots}
\usetikzlibrary{matrix,positioning,fit}

% Ruta a la carpeta de datos
\newcommand{\G}{datos_graficas}

\definecolor{cAzul}   {HTML}{1F4E79}
\definecolor{cNaranja}{HTML}{C55A11}
\definecolor{cVerde}  {HTML}{2E7D32}
\definecolor{cGris}   {HTML}{6E6E6E}
\definecolor{cRojo}   {HTML}{A61C1C}

% Estilo base: pensado para una columna IEEE
\pgfplotsset{
  estilo/.style={
    width            = 0.97\columnwidth,
    height           = 4.1cm,
    scale only axis  = false,
    grid             = major,
    grid style       = {gray!20, line width=0.25pt},
    tick label style = {font=\scriptsize},
    label style      = {font=\scriptsize},
    legend style     = {font=\scriptsize, draw=gray!45, fill=white,
                        fill opacity=0.85, text opacity=1,
                        inner sep=1.5pt, row sep=-1pt},
    legend cell align = left,
    every axis plot/.append style = {line width=0.7pt},
    scaled ticks = false,
    tick align   = outside,
    max space between ticks = 40pt,
  },
  ancho/.style = {width=0.47\textwidth, height=4.3cm},
  aligera/.style = {each nth point=#1, filter discard warning=false,
                    unbounded coords=discard},
  aligera/.default = 4,
}

% ---------- Codigo fuente ----------
\lstdefinestyle{estiloC}{
  language=C, basicstyle=\ttfamily\scriptsize,
  keywordstyle=\color{cAzul}\bfseries,
  commentstyle=\color{cGris}\itshape,
  frame=single, rulecolor=\color{gray!40},
  breaklines=true, showstringspaces=false, columns=fullflexible,
  xleftmargin=0.6em, aboveskip=4pt, belowskip=2pt
}
\lstset{style=estiloC}

% ---------- Marcas de trabajo pendiente ----------
% Cambie a \mostrarpendientesfalse antes de entregar
\newif\ifmostrarpendientes
\mostrarpendientesfalse
\newcommand{\pendiente}[1]{%
  \ifmostrarpendientes
    \par\noindent\fcolorbox{cRojo}{red!4}{%
      \parbox{0.95\columnwidth}{\scriptsize\color{cRojo}\textbf{PENDIENTE:} #1}}\par\smallskip
  \fi}

% ---------- Figura pendiente (marco vacio con la descripcion) ----------
\newcommand{\figpendiente}[2][3.2cm]{%
  \fbox{\parbox[c][#1][c]{0.92\columnwidth}{\centering\scriptsize\color{cRojo}%
  \textbf{FIGURA PENDIENTE}\\[2pt]#2}}}

% ---------- Atajos ----------
\newcommand{\xh}{\hat{x}}
\newcommand{\vh}{\hat{v}}
\newcommand{\thc}{\theta_{\mathrm{cmd}}}
\newcommand{\thlim}{\theta_{\mathrm{lim}}}
\newcommand{\Xm}{X_m}
\newcommand{\hef}{h_{\mathrm{ef}}}
\newcommand{\hsup}{h_{\mathrm{sup}}}
\newcommand{\Kp}{K_p}
\newcommand{\Kd}{K_d}
\newcommand{\Ki}{K_i}
\newcommand{\Ti}{T_i}
\newcommand{\Ts}{T_s}
